\documentclass{report}
\usepackage{amsmath,amssymb}
\setlength{\parindent}{0mm}
\setlength{\parskip}{1em}
\begin{document}
\begin{center}
\rule{6in}{1pt} \
{\large Xian-Ming Gu \\
{\bf Nested Krylov subspace methods based on the Hessenberg procedure for shifted linear systems}}

Institute of Mathematics and Computing Science \\ University of Groningen \\ Nijenborgh 9 \\ P O Box 407 \\ 9700 AK Groningen \\ The Netherlands
\\
{\tt guxianming@live.cn}\\
T.-Z.  Huang\\
B. Carpentieri\end{center}

In many computational applications it is required to solve simultaneously
shifted linear systems of the form
\begin{equation}
(A - \sigma_k I)x^{(k)} = b,\quad\ \sigma_k \in \mathbb{C},\quad k = 1,2,\ldots,s,
\label{eq1.1}
\end{equation}
where $A\in \mathbb{C}^{n\times n}$ is a non-Hermitian matrix, $x^{(k)}$
is the unknown vector and $b$ is the right-hand side, both of them are
complex vectors of length $n$. For instance, shifted linear systems arise
in the solution of PageRank problems and the rational approximation of
the action of a matrix function to a vector, i.e. $f(A)v$, where $v$ is
an arbitrary vector with length $n$.

In principle, a sequence of shifted linear systems (\ref{eq1.1}) can be
solved simultaneously at the cost of a single solve using shifted Krylov
subspace methods (abbreviated as KSMs). These methods employ the property
that Krylov subspaces are invariant under arbitrary diagonal shifts
$\sigma$ to the matrix $A$, i.e.,
\begin{equation}
\mathcal{K}_m(A, b) = \mathcal{K}_m(A - \sigma I, b),\quad \forall m \in
\mathbb{N}^+,\quad
\forall \sigma \in \mathbb{C}.
\label{eq1.2}
\end{equation}
As we know, the restarted GMRES method is a widely popular choice for
solving (\ref{eq1.1}), refer to \cite{AFUGR}. The computed GMRES shifted
residuals are not collinear in
general after the first restart so that it loses the computational
efficiency mentioned above. Consequently,
certain enforcement has to be made to guarantee the computed GMRES
shifted residuals collinear to each other in order to maintain the
computational efficiency. Note that in this case,
only the seed system satisfies the minimum residual property, the
solution of the other shifted systems is not equivalent to GMRES applied
to those systems, see \cite{AFUGR} for details. Moreover, another
attractive approach is to use the restarted FOM method for solving the
whole sequence
of shifted linear systems simultaneously, for all residuals are naturally
collinear \cite{AFUGR,VSRF}. As a result, the computational efficiency
can be maintained because the orthonormal basis and the Hessenberg matrix
are required to be calculated only once each time.

However, both the restarted GMRES method and the restarted FOM method for
(\ref{eq1.1}) are inherently derived by the Arnoldi procedure, which
tends to be expensive when $m$ becomes large in aspects of memory and
algorithmic cost. To alleviate these difficulties, our recent work
\cite{XMGTZH} proposed a new iterative strategy with collinear residuals
for solving sequence of shifted linear systems based on the Hessenberg
process, which is similar but cheaper compared to the Arnoldi procedure.
Our theoretical and numerical analyses indicate that iterative schemes
referred to as $\mathtt{sHessen(m)}$ may outperform the conventional
shifted FOM method in aspects of iteration counts and CPU time. From the
algorithmic property, it turns out that preserving orthonormality
in the construction of the Krylov basis is sometimes not necessary for an
efficient iterative solution of (\ref{eq1.1}). This observation motivates
us to extend the restarted Changing Minimal Residual method based on the
Hessenberg process (CMRH) \cite{HSCMRH}, which is similar
to the restarted shifted GMRES method but much cheaper, for solving (\ref{eq1.1}).

In practice, preconditioning shifted systems (\ref{eq1.1}) is necessary
to accelerate the convergence of shifted KSMs and this remains an open
question, because standard preconditioning techniques may destroy the
shift-invariance property (\ref{eq1.2}). In this work, we construct an
inner-outer scheme based on nested KSMs, where the inner solver is a
multi-shift KSM (i.e., $\mathtt{sHessen(m)}$) that acts as a
preconditioner for an outer flexible CMRH (FCMRH) method. Our
numerical experiments show that the new nested iterative scheme based on
the Hessenberg procedure is very competitive, and sometimes superior,
compared to the early established combination (FOM-FGMRES) in
\cite{MBMBG} in terms of CPU time. In summary, in our presentation first
we discuss the derivations of both the $\mathtt{sHessen(m)}$ method
\cite{XMGTZH} and the restarted shifted CMRH method. Then, we present
their flexible variant combined with nested inner-outer iterative
schemes, namely multi-shift Hessen-FCMRH.
Meanwhile, extensive numerical experiments involving PageRank problems
and numerical solutions of space(-time) fractional differential equations
are reported to illustrate the effectiveness of the restarted shifted
CMRH method and the Hessen-FCMRH method, respectively.

\begin{thebibliography}{99}
\bibitem{AFUGR}
A. Frommer, U. Gl\"{a}ssner, Restarted GMRES for shifted linear systems,
SIAM J. Sci. Comput., 19 (1998), pp. 15-26.

\bibitem{VSRF}
V. Simoncini, Restarted full orthogonalization method for shifted linear
systems, BIT., 43 (2003), pp. 459-466.

\bibitem{HSCMRH}
H. Sadok, CMRH: A new method for solving nonsymmetric linear systems
based on the Hessenberg reduction algorithm, Numer. Algorithms, 20
(1999), pp. 303-321.

\bibitem{MBMBG}
M. Baumann, M. B. van Gijzen, Nested Krylov methods for shifted linear
systems, SIAM J. Sci. Comput., 37 (2015), pp. S90-S112.

\bibitem{XMGTZH}
X.-M. Gu, T.-Z. Huang, G. Yin, B. Carpentieri, C. Wen, Restarted
Hessenberg method for shifted nonsymmetric linear systems with
applications to fractional differential equations, submitted to Numer.
Linear Algebra Appl., Aug. 20, 2015, 24 page. ({\it under review})

\end{thebibliography}


\end{document}
